%%% -*- coding: utf-8 -*-

\lecture{LINQ}{linq}
\section{LINQ}

\showtitlepage{Language Integrated Query}

\begin{frame}{LINQ (Language Integrated Query)}

  LINQ is a feature in C\# that builds query capabilities directly into the language syntax.

  \begin{itemize}
    \item Instead of writing custom loops and conditional statements to filter, sort, or transform data, LINQ allows you to write highly readable, declarative queries.
    \item The best part of LINQ is its versatility: you use the exact same syntax to query in-memory collections (like List<T> or arrays), databases (via Entity Framework), or XML files.
  \end{itemize}

\end{frame}

\begin{frame}{Collections: Java and C\#}
  \begin{center}
    \footnotesize
    \begin{tabular}{lll}
      \hline
      \textbf{Java} & \textbf{C\#} & \textbf{Kind}\\
      \hline
      \texttt{ArrayList<T>}  & \texttt{List<T>}            & class\\
      \texttt{List<T>}       & \texttt{IList<T>}           & interface\\
      \texttt{HashMap<K,V>}  & \texttt{Dictionary<K,V>}    & class\\
      \texttt{Map<K,V>}      & \texttt{IDictionary<K,V>}   & interface\\
      \texttt{HashSet<T>}    & \texttt{HashSet<T>}         & class\\
      \texttt{Set<T>}        & \texttt{ISet<T>}            & interface\\
      \texttt{Iterable<T>}   & \texttt{IEnumerable<T>}     & interface\\
      \hline
    \end{tabular}
  \end{center}
  \vspace{1ex}

  \hint{The trap for Java programmers: \texttt{List<T>} is a \emph{class} in
  C\#, corresponding to \texttt{ArrayList}, not an interface. Interfaces are
  spelled with a leading \texttt{I}, which is a naming convention rather than a
  language rule.}
\end{frame}

\begin{frame}{IEnumerable: The Common Denominator}
  \begin{itemize}
    \item \texttt{IEnumerable<T>} is what \texttt{foreach} needs and what every
      collection offers. It promises no more than that its elements can be
      visited one after the other.
    \item All LINQ operators are extension methods
      (\slideref{slide:extensions}) of \texttt{IEnumerable<T>}. That is why they
      are available on every collection without any collection knowing about
      LINQ.
    \item An \texttt{IEnumerable<T>} need not be a collection at all. It may
      compute its elements while being visited, which is what makes an infinite
      sequence or a query over a file possible.
  \end{itemize}
  \vspace{4ex}

  \hint{Declare parameters as \texttt{IEnumerable<T>} when you only iterate,
  and as \texttt{IList<T>} or \texttt{ISet<T>} when you need indexing, counting
  or membership tests.}
\end{frame}

\begin{frame}{The Operators we Actually Use in SEE}
  \begin{itemize}
    \item Filtering and mapping: \texttt{Where}, \texttt{Select},
      \texttt{SelectMany}, \texttt{Distinct}, \texttt{Cast},
      \texttt{OfType}.
    \item Asking questions: \texttt{Any}, \texttt{All}, \texttt{Contains},
      \texttt{Count}, \texttt{Sum}, \texttt{Min}, \texttt{Max}.
    \item Picking single elements: \texttt{First}, \texttt{FirstOrDefault},
      \texttt{Single}, \texttt{SingleOrDefault}.
    \item Sorting and grouping: \texttt{OrderBy}, \texttt{OrderByDescending},
      \texttt{ThenBy}, \texttt{GroupBy}.
    \item Materializing: \texttt{ToList}, \texttt{ToArray},
      \texttt{ToHashSet}, \texttt{ToDictionary}.
  \end{itemize}
  \vspace{4ex}

  \hint{\texttt{First} throws when nothing matches, \texttt{FirstOrDefault}
  yields \texttt{default(T)}, which is \texttt{null} for a reference type.
  \texttt{Single} additionally insists that there be no second match.}
\end{frame}

\begin{frame}[fragile]{Two Syntaxes, One Meaning}
  The examples use a class \texttt{Node} having the properties \texttt{Type} and
  \texttt{Id}, both of type \texttt{string}.
  \vspace{1ex}

  \begin{columns}[T]
    \column{0.5\textwidth}
      \textbf{Method syntax}
      \begin{lstlisting}[basicstyle=\scriptsize\ttfamily]
nodes.Where(n => n.Type == "Class")
     .OrderBy(n => n.Id)
     .Select(n => n.Id)
     .ToList();
      \end{lstlisting}
    \column{0.5\textwidth}
      \textbf{Query syntax}
      \begin{lstlisting}[basicstyle=\scriptsize\ttfamily]
(from n in nodes
 where n.Type == "Class"
 orderby n.Id
 select n.Id).ToList();
      \end{lstlisting}
  \end{columns}
  \vspace{1ex}

  The compiler translates query syntax into the very same calls, hence the two
  are interchangeable. Query syntax pays off for joins and for several
  \texttt{from} clauses, method syntax for short chains.
  \vspace{4ex}

  \hint{SEE currently uses method syntax exclusively: its sources contain no single query
  expression.}
\end{frame}

\begin{frame}[fragile]{LINQ in SEE}
  \begin{lstlisting}[basicstyle=\scriptsize\ttfamily]
// Graph.cs: all node types occurring in a graph
Nodes().Select(n => n.Type).ToHashSet();

// LSPImporter.cs: the first node isomorphic to a given one
graph.Nodes().FirstOrDefault(x => AreIsomorphic(node, x));

// SourceRangeIndex.cs: does every child satisfy a predicate?
files.Values.SelectMany(file => file.Children).All(IsHomomorphic);

// IncrementalTreeMap: an index of nodes by their ID
nodes.ToDictionary(n => n.ID, n => n);

// NodeSearch.cs: nodes grouped by their string representation
graph.Nodes().GroupBy(ElementToString)
             .ToDictionary(g => g.Key, g => g.ToList());
  \end{lstlisting}
  \vspace{1ex}

  About a fifth of SEE's C\# files use LINQ. Most frequent are
  \texttt{Contains}, \texttt{ToList}, \texttt{Select} and \texttt{Where}.
\end{frame}

\begin{frame}[fragile]{SEE Extends LINQ}
  As LINQ operators are nothing but extension methods, SEE adds its own, which
  then chain with the predefined ones:
  \vspace{1ex}

  \begin{lstlisting}[basicstyle=\scriptsize\ttfamily]
// DataModel/DG/GraphElementExtensions.cs
public static IEnumerable<T> WithGameObject<T>
    (this IEnumerable<T> elements) where T : GraphElement
  => elements.Where(HasGameObject);

// hence:
graph.Nodes().WithGameObject().Select(n => n.ID);
  \end{lstlisting}
  \vspace{4ex}

  \hint{\texttt{Utils/CollectionExtensions.cs} collects further operators of
  this kind, among them \texttt{GetOrAdd}, \texttt{RemoveWhere} and
  \texttt{Toggle}.}
\end{frame}

\begin{frame}[fragile]{Pitfalls}
  \begin{itemize}
    \item A query describes work rather than doing it. Enumerating it twice
      does the work twice, which is why an expensive query is usually closed
      with \texttt{ToList()}.
    \item Modifying the underlying collection while enumerating it throws,
      even when the query itself looks harmless.
  \end{itemize}
  \vspace{1ex}

  \begin{lstlisting}[language={},basicstyle=\scriptsize\ttfamily]
InvalidOperationException: Collection was modified;
enumeration operation may not execute.
  \end{lstlisting}
  \vspace{4ex}

  \hint{Materialize before you modify: iterate over
  \texttt{nodes.Where(\ldots).ToList()} when the loop body adds to or removes
  from \texttt{nodes}.}
\end{frame}

\begin{exerciseframe}{From Loop to Query}
  Replace the loop by a single LINQ chain yielding the same list.
  \vspace{1ex}

  \begin{lstlisting}[basicstyle=\scriptsize\ttfamily]
List<string> result = new List<string>();
foreach (Node node in nodes)
{
  if (node.Type == "Class")
  {
    result.Add(node.Id);
  }
}
  \end{lstlisting}
  \vspace{1ex}

  Which operator replaces the \texttt{if}, and which one the \texttt{Add}?
\end{exerciseframe}

\begin{solutionframe}{From Loop to Query}
  \begin{lstlisting}[basicstyle=\scriptsize\ttfamily]
List<string> result = nodes.Where(n => n.Type == "Class")
                           .Select(n => n.Id)
                           .ToList();
  \end{lstlisting}
  \begin{itemize}
    \item \texttt{Where} takes the role of the \texttt{if}, \texttt{Select}
      that of the expression handed to \texttt{Add}.
    \item \texttt{ToList} is what actually builds the list. Without it you get
      a description of the work, not a result.
  \end{itemize}
  \vspace{4ex}

  \hint{The chain says \emph{what} is wanted, the loop said \emph{how} to get
  it. That is the whole point of LINQ.}
\end{solutionframe}

\begin{exerciseframe}{Select or SelectMany?}
  A \texttt{Node} has a property \texttt{Children} of type
  \texttt{List<Node>}. Predict both numbers and, more importantly, both static
  types.
  \vspace{1ex}

  \begin{lstlisting}[basicstyle=\scriptsize\ttfamily]
// three nodes, carrying 2, 2 and 0 children
Console.WriteLine(nodes.Select(n => n.Children).Count());
Console.WriteLine(nodes.SelectMany(n => n.Children).Count());
  \end{lstlisting}
\end{exerciseframe}

\begin{solutionframe}{Select or SelectMany?}
  \begin{lstlisting}[language={},basicstyle=\scriptsize\ttfamily]
3
4
  \end{lstlisting}
  \begin{itemize}
    \item \texttt{Select} maps every node to its list, yielding an
      \texttt{IEnumerable<List<Node>>} of three lists.
    \item \texttt{SelectMany} concatenates those lists into a single
      \texttt{IEnumerable<Node>} of four children.
  \end{itemize}
  \vspace{1ex}

  \hint{Whenever a mapping yields a collection per element and you want the
  elements rather than the collections, \texttt{SelectMany} is the operator.
  SEE uses it to descend into children and file extensions.}
\end{solutionframe}

\begin{exerciseframe}{An Index by Type}
  Build a dictionary mapping every node type to the list of nodes of that type,
  in a single chain. What does it yield for the three nodes below?
  \vspace{1ex}

  \begin{lstlisting}[basicstyle=\scriptsize\ttfamily]
// A and B are of type "Class", C is of type "File".
Dictionary<string, List<Node>> byType = // your chain here
  \end{lstlisting}
  \vspace{1ex}

  Which operator groups, and which one turns the groups into a dictionary?
\end{exerciseframe}

\begin{solutionframe}{An Index by Type}
  \begin{lstlisting}[basicstyle=\scriptsize\ttfamily]
Dictionary<string, List<Node>> byType
  = nodes.GroupBy(n => n.Type)
         .ToDictionary(g => g.Key, g => g.ToList());

// Class = 2, File = 1
  \end{lstlisting}
  \begin{itemize}
    \item \texttt{GroupBy} yields a sequence of groups, each an
      \texttt{IEnumerable<Node>} carrying a \texttt{Key}.
    \item \texttt{ToDictionary} needs two functions: one for the key, one for
      the value.
  \end{itemize}
  \vspace{1ex}

  \hint{This is precisely what \texttt{NodeSearch} in SEE does to index all
  nodes of a graph. A duplicate key makes \texttt{ToDictionary} throw, hence
  group first.}
\end{solutionframe}

\begin{exerciseframe}{Your Own LINQ Operator}
  Write \texttt{WithType} as an extension method, so that it can be chained
  like any predefined operator.
  \vspace{1ex}

  \begin{lstlisting}[basicstyle=\scriptsize\ttfamily]
nodes.WithType("Class").Select(n => n.Id);   // A, B
  \end{lstlisting}
  \vspace{1ex}

  Where does it have to be declared, and what does it return so that
  \texttt{Select} can follow?
\end{exerciseframe}

\begin{solutionframe}{Your Own LINQ Operator}
  \begin{lstlisting}[basicstyle=\scriptsize\ttfamily]
public static class NodeExtensions
{
  public static IEnumerable<Node> WithType
      (this IEnumerable<Node> nodes, string type)
    => nodes.Where(n => n.Type == type);
}
  \end{lstlisting}
  \begin{itemize}
    \item A non-generic static class, with \texttt{this} on the first
      parameter, exactly as on \slideref{slide:extensions}.
    \item Returning \texttt{IEnumerable<Node>} is what keeps the chain open,
      and it keeps the operator lazy: nothing is filtered until enumerated.
  \end{itemize}
  \vspace{1ex}

  \hint{This is the shape of SEE's \texttt{WithGameObject}. Returning
  \texttt{List<Node>} instead would work too, but would enumerate eagerly and
  allocate.}
\end{solutionframe}